% ============================================================================
%  Problems: ตรรกศาสตร์พื้นฐาน (Basic Logic)
%  Ordered from easy (basic) → hard (POSN level).
%  Shared by exercise.tex and answer-key.tex
% ============================================================================

% --- Part 1: พื้นฐาน (Basic) ---
\examsection{ตอนที่ 1: พื้นฐาน — ประพจน์และตัวเชื่อม}

\problem{
  จงระบุว่าประโยคต่อไปนี้เป็นประพจน์หรือไม่ พร้อมระบุค่าความจริง (ถ้าเป็นประพจน์)
  \begin{enumerate}
    \item[(ก)] $2 + 2 = 5$
    \item[(ข)] ``เธอชื่ออะไร?''
    \item[(ค)] ประเทศไทยอยู่ในทวีปเอเชีย
    \item[(ง)] $x > 3$
  \end{enumerate}
}{
  (ก) เป็นประพจน์, ค่าความจริงเป็นเท็จ (F)

  (ข) ไม่เป็นประพจน์ (เป็นคำถาม)

  (ค) เป็นประพจน์, ค่าความจริงเป็นจริง (T)

  (ง) ไม่เป็นประพจน์ (มีตัวแปรอิสระ)
}

\problem{
  กำหนด $p$ แทน ``ฝนตก'' และ $q$ แทน ``ถนนเปียก''
  จงเขียนประพจน์ต่อไปนี้ในรูปสัญลักษณ์
  \begin{enumerate}
    \item[(ก)] ฝนตกและถนนเปียก
    \item[(ข)] ฝนไม่ตก
    \item[(ค)] ถ้าฝนตกแล้วถนนเปียก
    \item[(ง)] ฝนตกหรือถนนเปียก
  \end{enumerate}
}{
  (ก) $p \land q$

  (ข) $\lnot p$

  (ค) $p \to q$

  (ง) $p \lor q$
}

\problem{
  จงหาค่าความจริงของประพจน์ต่อไปนี้ เมื่อ $p = T$, $q = F$
  \begin{enumerate}
    \item[(ก)] $p \land q$
    \item[(ข)] $p \lor q$
    \item[(ค)] $\lnot p$
    \item[(ง)] $p \to q$
    \item[(จ)] $p \leftrightarrow q$
  \end{enumerate}
}{
  (ก) F, \quad (ข) T, \quad (ค) F, \quad (ง) F, \quad (จ) F
}

\problem{
  กำหนด $p$ = ``ฉันสอบผ่าน'' (T), $q$ = ``ฉันได้ไปค่าย'' (F)
  จงหาค่าความจริงของ $\lnot p \lor q$
}{
  $\lnot p = F$, \quad $F \lor F = F$

  ดังนั้น $\lnot p \lor q$ มีค่าความจริงเป็น \textbf{เท็จ}
}

% --- Part 2: การประยุกต์ (Intermediate) ---
\examsection{ตอนที่ 2: การประยุกต์ — ตารางค่าความจริงและการสมมูล}

\problem{
  จงสร้างตารางค่าความจริงของ $(p \land q) \to p$ และระบุว่าเป็นสัจนิรันดร์ (Tautology) หรือไม่
  (เป็นจริงทุกกรณี = Tautology)
}{
  \begin{tabular}{cc|c|c}
    $p$ & $q$ & $p \land q$ & $(p \land q) \to p$ \\
    \hline
    T & T & T & T \\
    T & F & F & T \\
    F & T & F & T \\
    F & F & F & T
  \end{tabular}

  เป็นจริงทุกกรณี $\rightarrow$ \textbf{เป็น Tautology}
}

\problem{
  จงตรวจสอบว่า $p \to q$ และ $\lnot q \to \lnot p$ สมมูลกันหรือไม่
  (สร้างตารางค่าความจริงเปรียบเทียบ)
}{
  \begin{tabular}{cc|c|c}
    $p$ & $q$ & $p \to q$ & $\lnot q \to \lnot p$ \\
    \hline
    T & T & T & T \\
    T & F & F & F \\
    F & T & T & T \\
    F & F & T & T
  \end{tabular}

  ค่าความจริงตรงกันทุกกรณี $\rightarrow$ \textbf{สมมูลกัน} (Contrapositive)
}

\problem{
  จงใช้กฎการสมมูลเปลี่ยนรูป $p \to q$ ให้อยู่ในรูปที่ไม่มีเครื่องหมาย $\to$
}{
  $p \to q \equiv \lnot p \lor q$
}

\problem{
  จงใช้นิเสธของตัวบ่งปริมาณเพื่อเขียนนิเสธของประโยคต่อไปนี้
  \begin{enumerate}
    \item[(ก)] $\forall x\; (x > 0)$ \quad (ทุกจำนวนมีค่ามากกว่า 0)
    \item[(ข)] $\exists x\; (x^2 = 9)$ \quad (มีจำนวนที่ยกกำลังสองแล้วได้ 9)
  \end{enumerate}
}{
  (ก) $\lnot \forall x\; (x > 0) \equiv \exists x\; (x \leq 0)$

  (ข) $\lnot \exists x\; (x^2 = 9) \equiv \forall x\; (x^2 \neq 9)$
}

\problem{
  จงเขียนนิเสธของประโยค ``นักเรียนทุกคนชอบคณิตศาสตร์''
  โดยใช้ตัวบ่งปริมาณ
}{
  ให้ $P(x)$ แทน ``$x$ ชอบคณิตศาสตร์''

  ประโยคเดิม: $\forall x\; P(x)$

  นิเสธ: $\exists x\; \lnot P(x)$

  ``มีนักเรียนอย่างน้อย 1 คนที่ไม่ชอบคณิตศาสตร์''
}

\problem{
  จงใช้ Modus Ponens หรือ Modus Tollens ในการสรุปผล
  \begin{enumerate}
    \item[(ก)] ถ้าฝนตก แล้วฉันจะไม่ไปโรงเรียน \quad ฝนตก \quad ดังนั้น?
    \item[(ข)] ถ้าเธอเป็นนก แล้วเธอต้องมีปีก \quad เธอไม่มีปีก \quad ดังนั้น?
  \end{enumerate}
}{
  (ก) Modus Ponens: ฉันจะไม่ไปโรงเรียน

  (ข) Modus Tollens: เธอไม่ใช่นก
}

\problem{
  จงระบุว่าการอ้างเหตุผลต่อไปนี้สมเหตุสมผลหรือไม่
  \begin{enumerate}
    \item[(ก)] ถ้าสอบได้คะแนนดี แล้วจะได้เกรด A \quad ได้เกรด A \quad ดังนั้น สอบได้คะแนนดี
    \item[(ข)] ถ้าเป็นคนขยัน แล้วจะสอบผ่าน \quad ไม่ขยัน \quad ดังนั้น สอบไม่ผ่าน
  \end{enumerate}
}{
  (ก) \textbf{ไม่สมเหตุสมผล} (Converse Error — อาจได้ A เพราะเหตุผลอื่น)

  (ข) \textbf{ไม่สมเหตุสมผล} (Inverse Error — อาจสอบผ่านเพราะโชคช่วย)
}

\problem{
  กำหนดประพจน์ $\lnot(p \lor \lnot q)$
  จงใช้กฎเดอมอร์แกนและกฎการสมมูล เปลี่ยนให้อยู่ในรูปอย่างง่าย
}{
  $\lnot(p \lor \lnot q) \equiv \lnot p \land \lnot(\lnot q)
  \equiv \lnot p \land q$
}

\problem{
  มีคน 2 คนคือ A และ B โดยแต่ละคนเป็นคนพูดจริง (T) หรือคนพูดเท็จ (L)

  A กล่าวว่า ``B เป็นคนพูดจริง''

  B กล่าวว่า ``พวกเราทั้งคู่เป็นคนละประเภทกัน''

  จงหาว่า A และ B เป็นคนประเภทใด
}{
  \textbf{กรณี 1: A พูดจริง} $\rightarrow$ B พูดจริง $\rightarrow$ B พูดว่า ``A กับ B เป็นคนละประเภทกัน'' เป็นจริง $\rightarrow$ ขัดแย้ง (เพราะทั้งคู่เป็น T เหมือนกัน)

  \textbf{กรณี 2: A พูดเท็จ} $\rightarrow$ B พูดเท็จ $\rightarrow$ B พูดว่า ``A กับ B เป็นคนละประเภทกัน'' เป็นเท็จ $\rightarrow$ ทั้งคู่เป็นประเภทเดียวกัน (L ทั้งคู่) $\checkmark$

  \textbf{สรุป:} A พูดเท็จ, B พูดเท็จ
}

% --- Part 3: ระดับ สอวน. (POSN Level) ---
\examsection{ตอนที่ 3: ระดับข้อสอบ สอวน.}

\problem{
  มีคน 3 คนคือ A, B และ C โดยแต่ละคนพูดจริงหรือพูดเท็จ

  A กล่าวว่า ``B เป็นคนพูดเท็จ''

  B กล่าวว่า ``A และ C เป็นคนประเภทเดียวกัน''

  C กล่าวว่า ``A เป็นคนพูดจริง''

  จงหาว่าใครเป็นคนประเภทใด
}{
  \textbf{กรณี 1: A พูดจริง}
  $\rightarrow$ B พูดเท็จ (เพราะ A กล่าวว่า ``B เป็นคนพูดเท็จ'')

  B พูดเท็จ $\rightarrow$ ``A และ C ประเภทเดียวกัน'' เป็นเท็จ $\rightarrow$ A กับ C คนละประเภท $\rightarrow$ C พูดเท็จ

  แต่ C พูดว่า ``A พูดจริง'' ซึ่งเป็นจริง $\rightarrow$ \textbf{ขัดแย้ง!} (C พูดเท็จ แต่ดันพูดความจริง)

  \medskip
  \textbf{กรณี 2: A พูดเท็จ}
  $\rightarrow$ B พูดจริง (เพราะ A โกหกว่า ``B เป็นคนพูดเท็จ'')

  B พูดจริง $\rightarrow$ ``A และ C ประเภทเดียวกัน'' เป็นจริง $\rightarrow$ C พูดเท็จเหมือน A

  ตรวจสอบ C: C พูดว่า ``A พูดจริง'' — C พูดเท็จ และประโยคนี้เป็นเท็จ (เพราะ A พูดเท็จ) $\checkmark$ สอดคล้อง

  \medskip
  \textbf{สรุป:} A พูดเท็จ, B พูดจริง, C พูดเท็จ
}

\problem{
  (ข้อสอบจริงปี 65 ข้อ 1) กำหนด "คอนทราดิกชัน" คือประพจน์ที่มีค่าความจริงเป็นเท็จทุกกรณี ประพจน์ในข้อใดเป็นคอนทราดิกชัน
  \begin{enumerate}
    \item[ก.] $(p \land q)\land \lnot(p \lor q)$
    \item[ข.] $(p \lor q)\land (\lnot p \lor q)$
    \item[ค.] $(p \rightarrow q)\land \lnot(p \lor q)$
    \item[ง.] $(p \leftrightarrow  q)\land (p \lor q)$
  \end{enumerate}
}{
  ก.
}

\problem{
  (ข้อสอบจริงปี 65 ข้อ 2) ผลในข้อใดที่ทำให้การอ้างเหตุผลนี้ สมเหตุสมผล \newline
  เหตู: 
  \begin{enumerate}
    \item ถ้าฉันขยันแล้วฉันทำข้อสอบได้
    \item ถ้าฉันไม่ขยันแล้วพ่อแม่เสียใจ
    \item ถ้าฉันเรียนในมหาวิทยาลัยแล้วพ่อแม่ไม่เสียใจ
    \item ฉันทำข้อสอบไม่ได้
  \end{enumerate}
  -----------------------------------------
  \begin{enumerate}
    \item[ก.] ฉันขยัน
    \item[ข.] ฉันทำข้อสอบได้
    \item[ค.] พ่อแม่เสียใจ
    \item[ง.] ฉันเรียนในมหาวิทยาลัย
  \end{enumerate}
}{
  ก.
}

\problem{
  (ข้อสอบจริงปี 68 ข้อ 18) กำหนดบัตร 5 ใบ วางบนโต๊ะ หงายหน้าที่มีเขียนว่า P, T, 5, 6 และ 8 บัตรแต่ละใบด้านนึงเป็นตัวเลข ด้านนึงเป็นตัวอักษร
  ถ้าต้องการพิสูจน์ว่า "ถ้าด้านนึงเป็นสระ (ภาษาอังกฤษ) แล้วอีกด้านนึงเป็นจำนวนเต็มบวกคู่" ต้องพลิกบัตรที่เขียนว่าอะไร
}{
  5
}

\problem{
  (ข้อสอบจริงปี 68 ข้อ 48) มีกล่อง 3 ใบ ข้างในกล่องมี เงิน หรือ ทอง อย่างใดอย่างหนึ่ง และมีการติดป้ายหน้าว่า "ทอง" "เงิน" และ "ทองหรือเงิน" แต่มีเพียงกล่องเดียวที่ติดป้ายตรงกับสิ่งของภายในกล่อง ถ้าคุณเปิดกล่องที่ติดป้ายว่าทองหรือเงินแล้วพบว่าเป็นทอง จงหาว่ากล่องที่มีเงินอยู่ข้างในคือกล่องที่ติดป้ายใด
}{
  ทอง
}